\section{Introduction}\label{sec:introduction}

Directional transience says that a random walk eventually leaves every
backward half-space. Ballisticity asks for more: linear progress with a
strictly positive speed. For a random walk in a random environment,
repeated visits to unfavorable regions can separate these two behaviors.
We prove that they coincide for nearest-neighbor walks in iid uniformly
elliptic environments in every dimension at least two.

\subsection{Model and main result}

Let $d\ge2$, let $e_1,\ldots,e_d$ be the coordinate unit vectors, and set
\[
 U=\{\pm e_1,\ldots,\pm e_d\},\qquad
 \Delta_U=\left\{p\in[0,1]^U:\sum_{e\in U}p(e)=1\right\}.
\]
Let $\nu$ be a probability measure on $\Delta_U$. An environment is a
collection $\omega=(\omega(x,\cdot))_{x\in\mathbb Z^d}$ of transition
rows, with product law
\[
 P=\nu^{\otimes\mathbb Z^d}.
\]
Thus different sites carry independent, identically distributed rows;
no independence among the entries of one row is required. Throughout,
we assume \emph{uniform ellipticity}: there is a constant $\kappa>0$ such
that
\[
 \nu\!\left(\min_{e\in U}p(e)\ge\kappa\right)=1.
\]
For fixed $\omega$, the quenched law $P_{x,\omega}$ is that of the Markov
chain starting at $x$ and satisfying
\[
 P_{x,\omega}(X_{n+1}=y+e\mid X_n=y)=\omega(y,e).
\]
Its annealed law is $P_x(\cdot)=\int P_{x,\omega}(\cdot)P(d\omega)$.
We also use $P_x$ for the joint environment-path law when both coordinates
are needed. Write $E$ for environment expectation, and use the corresponding
law labels on other expectations when necessary.

For a fixed $\ell\in S^{d-1}$, define
\[
 A_\ell=\{X_n\cdot\ell\longrightarrow+\infty\}.
\]
The walk is \emph{directionally transient} in direction $\ell$ if
$P_0(A_\ell)=1$, and \emph{ballistic} in that direction if
\[
 P_0\!\left(\liminf_{n\to\infty}\frac{X_n\cdot\ell}{n}>0\right)=1.
\]

\begin{theorem}[Directional transience implies ballisticity]\label{thm:main}
Let $d\ge2$ and let $P=\nu^{\otimes\mathbb Z^d}$ be an iid law of
nearest-neighbor transition rows satisfying uniform ellipticity. If
$\ell\in S^{d-1}$ and $P_0(A_\ell)=1$, then there exists a deterministic
vector $v=v(\nu)\in\mathbb R^d$ such that
\[
 v\cdot\ell>0,
 \qquad
 P_0\!\left(\lim_{n\to\infty}\frac{X_n}{n}=v\right)=1.
\]
In particular, the walk is ballistic in the same direction $\ell$.
\end{theorem}

The direction may be any fixed real unit vector. No drift, symmetry,
reversibility, independence within a row, or regeneration-time moment
assumption is added. The velocity is a property of the row law, since an
almost-sure deterministic vector limit is unique. The assertion is for
each fixed direction and does not require a common exceptional null set
for all directions.

\subsection{Context and significance}

Theorem~\ref{thm:main} resolves positively the ballisticity conjecture for
iid uniformly elliptic nearest-neighbor environments in dimension at
least two. Fribergh and Kious~\cite[Conjecture~1.1]{FriberghKious2016}
formulate it with precisely the one-direction transience assumption used
here. The distinction between escape and speed is essential: already in
one dimension, uniformly elliptic iid environments can give a transient
walk with zero speed; see Solomon~\cite{Solomon1975} and the discussion in
\cite[p.~510]{Sznitman2002}. The higher-dimensional conjecture asks whether
the additional spatial routes prevent this form of slowdown.

The regeneration construction of Sznitman and
Zerner~\cite{SznitmanZerner1999} separates the walk at new height records
below which it never returns. After an initial piece, the resulting
pieces are independent with a common conditioned law; see also
\cite[Section~1]{Sznitman2002} and
\cite[Section~2]{FriberghKious2016}. A finite mean duration then gives a
velocity by the renewal law of large numbers. Almost-sure finiteness of
the pieces does not supply that mean. The proof here first controls a
piece's spatial radius---its greatest distance from its starting point---
and only later its duration. Zeitouni's
survey~\cite{Zeitouni2006} gives further background on this distinction
and the regeneration method.

Several quantitative routes to ballisticity precede the present result.
Kalikow's auxiliary-walk approach~\cite{Kalikow1981} led to sufficient
directional drift conditions. Under Kalikow's condition,
Sznitman~\cite{Sznitman2000} obtained slowdown estimates and a central
limit theorem through regeneration-time tails. His conditions
$(T)_\gamma$ and $(T')$ require stretched-exponential control of backward
slab exits in nearby directions, and his effective criterion gives a
finite-box characterization of $(T')$ yielding positive
speed~\cite{Sznitman2002}. Berger, Drewitz, and
Ram\'irez~\cite{BergerDrewitzRamirez2014} proved that sufficiently strong
polynomial exit bounds imply $(T')$. Guerra and
Ram\'irez~\cite{GuerraRamirez2020} proved $(T')\Rightarrow(T)$, where
$(T)$ is the exponential slab-exit condition. Guerra's later
preprint~\cite[Definition~1.4 and Theorem~1.9]{Guerra2020WeakCriterion}
gives a finite-box criterion already implied by backward-exit
probabilities of order $o(L^{1-d})$ in nearby directions. These results
start from quantitative exit assumptions. Theorem~\ref{thm:main}
starts with almost-sure escape in one direction and derives the
quantitative crossing estimate needed for finite mean duration.

Even the spatial intermediate conclusion requires care.
Simenhaus~\cite[Theorem~1]{Simenhaus2007} obtains a deterministic
asymptotic direction from transience in a nonempty open set of
directions. Our spatial-radius estimate uses the specified single
direction and gives a limiting ray before a duration moment has been
proved. The record-counting proof is related to the cone-renewal
argument in \cite[Lemma~2]{Simenhaus2007}; its treatment of arbitrary
real projected heights is given directly in Section~\ref{sec:geometry}.

The two-walk construction has a close precedent in the common spatial
regeneration levels of Rassoul-Agha and
Sepp\"al\"ainen~\cite[Section~7]{RassoulAghaSeppalainen2009}. Comparing
shared and independent copies to control quenched fluctuations belongs
to the second-moment approach of Bolthausen and
Sznitman~\cite{BolthausenSznitman2002} and Berger and
Zeitouni~\cite{BergerZeitouni2008}. Their invariance principles assume
additional ballisticity or regeneration-time moments. Here the
comparisons use spatial first moments, truncated variances, and selected
Gaussian scales; the quenched approximation is proved along those
scales in environment probability.

Finally, normalized successful-endpoint distributions can spread out or
split into groups whose mutual distances diverge. Translation-invariant
compactifications retaining these possibilities were developed by
Mukherjee and Varadhan~\cite{MukherjeeVaradhan2016} and, for polymer
endpoints, by Bates and Chatterjee~\cite{BatesChatterjee2020}. We construct
the profile and upper-environment array needed here and prove its
finite-window entropy bound. That bound excludes diffuse mass and
infinitely many separated groups on the event of persistent survival
loss. An order comparison in the two distant tails of a remaining
profile then contradicts that loss.

\subsection{Proof strategy}

The proof separates spatial control from control of elapsed time.
Section~\ref{sec:geometry} constructs true record words in the original
real direction. Counting ordinary record indices gives finite mean
projected width, even when projected heights are not discrete. A forward
template at a tilted record then gives finite mean spatial radius.
This produces a deterministic limiting ray without a duration moment and
permits a signed coordinate permutation so that height is $x_1$.

For that coordinate height, let
\[
 a_N=P_{0,\omega}(T_N<T_{-1})
\]
be the quenched probability of reaching level $N$ before level $-1$.
Here $T_j$ denotes the first hit of level $j$. The decisive estimate is
\eqref{eq:5}: for a fixed $0<\beta<1$, the environment probability of
$a_N<N^{-\beta}$ decays faster than every power of $N$. The same section
proves that this estimate supplies the missing duration moment and hence
Theorem~\ref{thm:main}. The rest of the paper proves the estimate.

Sections~\ref{sec:fluctuations} and \ref{sec:gaussian} compare two walks
at first hits of integer layers. Independent common regeneration words
give symmetric increments and a scale defined by their truncated variance.
On Gaussian subsequences, a shared-environment comparison must exclude
positive limiting occupation on the projected diagonal. Its cutoff-scale
argument also controls the smaller scales that need not be Gaussian.
The resulting two-copy limit yields a quenched approximation without
assuming a second moment or a pre-existing quenched invariance principle.

Section~\ref{sec:clipping} combines Gaussian and non-Gaussian subscales
to retain positive quenched mass whose first-hit positions lie in narrow
tubes at every large scale. Section~\ref{sec:kernels} turns this into separated endpoint seeds,
rapid-loss bounds, and an outward-order estimate with an exponentially
integrable local remainder. All adaptive tests use only departure rows
below their tested top layer.

If \eqref{eq:5} fails, conditioning on its bad environments costs only
logarithmic relative entropy. The multiscale stages of
Section~\ref{sec:stages} then give a positive expected loss of global survival mass on episodes
with bounded stage count. Each episode is an adaptively chosen stretch
of layers; the bound on its stage count also controls the logarithm of
its height. In Section~\ref{sec:stationary}, averaging
over episodes gives a stationary law of local endpoint profiles and
their upper environments. Entropy bounds exclude diffuse mass and
infinitely many separated profile classes on the positive-drop event.
A tail-order statistic in a surviving full-support class contradicts
the sustained drop: testing far out in its two tails removes the local
environmental bias from the outward-order remainder. This contradiction
proves \eqref{eq:5} and completes the speed argument.

\subsection{Positive-probability transience}

In dimensions at least three, a separate directional zero--one theorem
allows us to assume only positive transience probability. The input is
Theorem~1.1 of the companion manuscript~\cite{OpenAIDirectionalZeroOne2026}. Together with Theorem~\ref{thm:main}, it also identifies
exactly the directions in which escape has positive probability.

\begin{corollary}[Positive-probability transience and the velocity hemisphere]\label{cor:positive-transience}
Let $d\ge3$ and let $P=\nu^{\otimes\mathbb Z^d}$ be an iid law of
nearest-neighbor transition rows satisfying uniform ellipticity. For a
fixed $\ell\in\mathbb R^d\setminus\{0\}$, set
$A_\ell=\{X_n\cdot\ell\longrightarrow+\infty\}$. If $P_0(A_\ell)>0$,
then there is a deterministic nonzero vector $v=v(\nu)$ such that
\[
 v\cdot\ell>0,
 \qquad
 P_0\!\left(\lim_{n\to\infty}\frac{X_n}{n}=v\right)=1.
\]
If such a direction exists, then the deterministic sets of directions
satisfy
\[
 \{u\in S^{d-1}:P_0(A_u)=1\}
 =\{u\in S^{d-1}:P_0(A_u)>0\}
 =\{u\in S^{d-1}:v\cdot u>0\}.
\]
In particular, $P_0(A_u)=0$ for each fixed $u\in S^{d-1}$ with
$v\cdot u\le0$, including every equator direction. The equality concerns
annealed probabilities for each fixed deterministic direction, not a
simultaneous pathwise assertion over all directions.
\end{corollary}

\begin{proof}
Uniform ellipticity implies the strict ellipticity assumed by the
companion's directional zero--one law. It therefore upgrades
$P_0(A_\ell)>0$ to $P_0(A_\ell)=1$. Since
$A_\ell=A_{\ell/\|\ell\|}$, Theorem~\ref{thm:main} gives the stated
velocity and positive projection.

Fix this vector $v$ and a direction $u\in S^{d-1}$. If $v\cdot u>0$,
the vector law of large numbers gives $X_n\cdot u/n\to v\cdot u>0$
almost surely, hence $P_0(A_u)=1$. Conversely, if $P_0(A_u)>0$, the
preceding argument applied to $u$ gives a deterministic vector $w$ with
$X_n/n\to w$ almost surely and $w\cdot u>0$. The two probability-one
convergence events intersect, and uniqueness of a vector limit gives
$w=v$. Thus $v\cdot u>0$. These two implications prove the equality of
direction sets and the assertion on the closed opposite hemisphere.
\end{proof}

\paragraph{Conventions.}
Logs are natural. Positive constants may change from line to line and
may depend on the fixed row law and dimension; additional dependencies
are indicated when uniformity matters. Walks in a shared environment
are conditionally independent given that environment. At a finite path
prefix, only rows at departure sites have contributed transition factors.
This distinction is retained when a path stops on first arrival at a
layer or a previously visited set.
